\clearpage\thispagestyle{empty}
\pdfbookmark[0]{Quelle zum Brief an von Lindenau vom 30. August 1808}{lindenau1808aug30-sources}
\begin{center}\textbf{Quelle zum Brief an von Lindenau vom 30. August 1808}\end{center}
\noindent Diese KI-Übertragung ins heutige Deutsch folgt dem Brief von Carl Friedrich Gauß an Bernhard August von Lindenau vom 30. August 1808 in Gauß’ \href{https://gdz.sub.uni-goettingen.de/id/PPN236007467}{\textit{Werke}, Bd.~VI} (Göttingen, 1874), S.~307. Der Erstdruck erschien als Nr.~XXXIV, \textit{Auszug aus einem Schreiben des Herrn Professors Gauß}, im Septemberheft der \textit{Monatlichen Correspondenz}, Bd.~XVIII, S.~\href{https://www.digitale-sammlungen.de/en/view/bsb10538609?page=291}{269}, \href{https://www.digitale-sammlungen.de/en/view/bsb10538609?page=292}{270}, \href{https://www.digitale-sammlungen.de/en/view/bsb10538609?page=293}{271}, \href{https://www.digitale-sammlungen.de/en/view/bsb10538609?page=294}{272} und \href{https://www.digitale-sammlungen.de/en/view/bsb10538609?page=295}{273}.
\par\medskip
Der Erstdruck wiederholt hier die Bahnelemente. Die Werkausgabe verweist dafür auf die eigene Juno-Mitteilung und stellt die Ausführungen zur Polhöhe hinter den zugehörigen Aufsatz auf S.~146--147. Eine englische Teilübersetzung findet sich bei \href{https://doi.org/10.1007/978-3-319-58118-7_12}{Clifford J. Cunningham, \textit{Scientific Papers} (Springer, 2017), S.~283--284}. Diese Übertragung folgt der Gliederung und den Tabellenwerten der Werkausgabe. KI-Übertragung und KI-Satz.
\clearpage
